\section{Theoretical Introduction}
Gaussian elimination with partial pivoting consists of 3 main steps:
\subsection{Transform matrix into upper-triangular matrix}

\subsubsection{Starting conditions}
We start with the system of linear equations looking like this:

\[
\begin{matrix}

&a_{11}x_1 &{}+&a_{12}x_2&+&\dots&+&a_{1n}x_n &=&b_1,\\

&a_{21}x_1 &{}+&a_{22}x_2&+&\dots&+&a_{2n}x_n &=&b_2,\\

&\vdots    &&\vdots      & &     & &  \vdots  & &\vdots\\

&a_{n1}x_1&{}+&a_{n2}x_2&+&\dots &+&a_{nn}x_n&=&b_n.

\end{matrix}
\]

In order for this method to work all the elements of "diagonal" line:
\[ a_{11}, a_{22}, \dots, a_{nn} \]
Must be different from zero since We will be dividing by them.

We will denote rows as '$w_i$' where 'i' is number of the row.

\subsubsection{Zeroing first column}
We start transforming the system by "zeroing" elements in first column excluding first row element. We do it by multiplying first row by $l_{i1}$, where:
\[ l_{i1} = \frac{ a_{i1}^{(1)} }
{ a_{11}^{(1)} }  \]
And then substracting what we got ($ l_{i1}w_1 $), from \textit{i} row.

Doing so we obtain a system of linear equations:

\[
\begin{matrix}

&a_{11}x_1 &{}+&a_{12}x_2&+&\dots&+&a_{1n}x_n &=&b_1,\\

&0 &{}+&(a_{22} - a_{12}l_{21})x_2&{}+&\dots&{}+&(a_{2n} - a_{1n}l_{21})x_n &=&b_2 - b_{1}l_{21},\\

&\vdots    &&\vdots      & &     & &  \vdots  & &\vdots\\

&0&{}+&(a_{n2} - a_{12}l_{n1})x_2&+&\dots &+&(a_{nn} - a_{1n}l_{n1})x_n&=&b_n - b_{1}l_{n1}.

\end{matrix}
\]

\subsubsection{Zeroing second column}
We continue onto the second column, this time we will zero all elements except first and second rows.
Row multiplier becomes:
\[ l_{i2} = \frac{ a_{i2}^{(2)} }{ a_{22}^{(2)} } \]
Where:
\[ a_{22}^{(2)} = (a_{22} - a_{12}l_{21}) \]
And:
\[ a_{i2}^{(2)} = (a_{i2} - a_{12}l_{i1}) \]
\newpage
They are modified values obtained from previous step.
We continue as in the first step and we end up with:

\[
\begin{matrix}

&a_{11}x_1 &{}+&a_{12}x_2&+&\dots&+&a_{1n}x_n &=&b_1,\\

&0 &{}+& a_{22}^{(2)}x_2&{}+&\dots&{}+& a_{2n}^{(2)}x_n &=&b_2^{(2)},\\

&\vdots    &&\vdots      & &     & &  \vdots  & &\vdots\\

&0 &{}+& 0 &{}+&\dots&{}+& a_{nn}^{(3)}x_n &=&b_2^{(3)},

\end{matrix}
\]

\subsubsection{Zeroing next columns}
We repeat this process $n-1$ times and we end up with upper triangular matrix:

\[
\begin{matrix}

&a_{11}x_1 &{}+&a_{12}x_2&+&\dots&+&a_{1n}x_n &=&b_1,\\

&0 &{}+& a_{22}^{(2)}x_2&{}+&\dots&{}+& a_{i2}^{(2)}x_n &=&b_2^{(2)},\\

&\vdots    &&\vdots      & &     & &  \vdots  & &\vdots\\

&0 &{}+& 0 &{}+&\dots&{}+& a_{nn}^{(n)}x_n &=&b_2^{(n)},

\end{matrix}
\]

\subsection{Backward substitution}
After transforming the system we solve the system from last to first. \\
First we calculate value of last element:
\[ x_n = \frac{b_n}{a_{nn}} \]
Then one "above":
\[ x_{n-1} = \frac{ b_{n-1} - a_{n-1, n}x_n}{a_{n-1, n-1}} \]
And so on, for $x_k$:
\[ x_{k} = \frac{b_k - \sum_{j = k + 1}^n a_{kj}x_j}{a_{kk}} \]

\newpage
\subsection{Partial Pivoting}
Gaussian elimination method has one flaw, where it can come into halt if:
\[ a_{kk}^{(k)} = 0 \]
To avoid it we use method of pivoting, in our case we will use partial pivoting method.
We do it before each Gaussian elimination step since this will lead to smaller error.
\\We first find a row $i$ such that:
\[ |{a_{ik}^{k}}| = \underset{j}{max} \{ |{a_{kk}^{k}}|, |{a_{k+1, k}^{k}}|, \cdots, |{a_{nk}^{k}}|\} \]

Then we swap this row with k-th row. Since the matrix we use is assumed to be nonsingular then $|{a_{ik}^{k}}| \neq 0$ will be always true. After that we continue with the Gaussian elimination method.

Let's compare full pivoting to partial pivoting.
Full pivoting carries more computational load. It comes from the fact that:
In full pivoting we need to compare absolute values of the matrix elements. (Which gives us $ k^2-1$ comparisons every step as opposed to $k-1$ for partial pivoting). We also have column interchanges and corresponsing interchanges in the order of elements of \textbf{x}.
Therefore when it comes to speed, partial pivoting is faster than full pivoting.

\newpage
